\clearpage
\section{Local route generation, budget sensitivity and measurement power}
\label{sec:review_followups}

\subsection{Local inhibitory gating in the conductance task}
\label{sec:local_conductance_gate}
This study was specified after exposure to a referee panel's exploratory gate results. Those runs are excluded here. The protocol and numerical adapter were committed before execution on twenty new seed blocks, 2026090801--2026090820. Within each seed, aligned and opposed targets and all nine rules shared teacher perturbations, initial students, features, contexts and minibatch streams. The model, nominal gains and independent log perturbations are defined in the main Methods. Each condition represents sixteen terminal E/I gains, two proximal inhibitory gains, four daughter couplings and two somatic couplings.

Let $\delta=(V_s-y)/\operatorname{Var}_{\rm train}(y)$ denote the somatic error and $a_p\in\{0,10\}$ the inhibitory activity at parent $p$. The hard local rule multiplies terminal eligibilities by $\delta\bm1[a_p=0]$, while proximal and somatic eligibilities receive $\delta$. It uses neither exact paths nor exact-field projections during learning. A secondary continuous gate uses $\delta/(1+g_p^{\rm I}a_p)$ at terminals, again retaining unit proximal and somatic multipliers. Both receive an imposed inhibitory context and an output error; neither generates those signals autonomously.

Exact, uniform, calibrated and three-profile oracle rules retain the previous comparators. Calibration averages initial exact path derivatives on 512 unlabeled examples. Two weighted leaf-subtree profiles with oracle projection amplitudes and unit proximal credit test the context distinction separately from proximal provision. The earlier three-profile oracle also includes a weighted proximal-pair column. The new control therefore has two distal profiles plus a fixed proximal rule, rather than an entirely two-dimensional six-site dictionary. Further controls reverse the distal gate or apply it at the proximal compartments as well. In the latter, inhibitory eligibility is nonzero only when inhibition is present, while its credit multiplier is then zero; both proximal inhibitory gradients vanish identically.

Each task supplies 2,048 training, 1,024 validation and 4,096 test examples without added label noise. Adam uses batches of 128, $\beta_1=0.9$, $\beta_2=0.999$, $\epsilon=10^{-8}$, no weight decay, norm clipping at ten and log-conductance bounds $[-7,7]$. The primary rate 0.03 was retained from the earlier protocol; 0.01 and 0.1 were symmetric robustness conditions for every rule. All $20\times2\times9\times3=1,080$ trajectories continued uninterrupted to 16,384 updates. Minimum validation NMSE selected a checkpoint within the primary 4,096-update window and the extended window. Fixed endpoints and every saved curve point are retained separately.

Two prespecified paired contrasts used exact two-sided sign-flip tests: opposed-task broadcast-minus-local-gate NMSE, and the difference in this gap between opposed and aligned tasks. Both were positive in all twenty seeds, with Holm-adjusted $P=3.81\times10^{-6}$ in the two-test family. Intervals use 20,000 whole-seed bootstrap draws. Other rates, budgets and contrasts are descriptive. The local-minus-exact opposed-task difference was $1.49\times10^{-6}$ (95\% interval, $-1.60\times10^{-6}$ to $4.86\times10^{-6}$), and the maximum local-gate seed error was $5.47\times10^{-5}$. These satisfy the frozen practical precision reference without establishing asymptotic equivalence.

At the primary rate, exact, proper local gates and both oracle dictionaries contacted no conductance bounds through 16,384 updates; eighteen of twenty opposed broadcast fits had contacted bounds by 4,096 and all twenty by 16,384. No fit activated gradient clipping. The proper local gate changed all 24 parameters, whereas proximal gating froze both inhibitory gains. Ten mechanism tests verify derivatives, absence of exact-path access in the local update code, and these eligibility interactions. An excluded historical-seed replay reproduces earlier exact, unit and calibrated trajectories bit for bit; the independently projected dictionary agrees to floating-point association error. All protocols, controls, states and exclusions accompany Source Data.

\begin{figure}[tbp]
\centering
\includegraphics[width=\textwidth]{supplementary/figure_S53_local_gate_controls.pdf}
\caption{\textbf{Local conductance-gate controls across rates and budgets.} \textbf{A,B}, Aligned targets within 4,096 and 16,384 updates. \textbf{C,D}, Opposed targets within the same windows. Entries are mean test NMSE at validation-selected states across twenty new paired seeds; color shows its base-ten logarithm. Every rule uses every rate; 0.03 is primary. \textbf{E}, Separate historical twenty-seed cohort (2101--2120). At identical initial or broadcast-trained weights, cancellation is the norm of mixed context-conditional distal gradients divided by the mixture of their norms; zero means complete cancellation. Symbols show means and bars descriptive 95\% whole-seed bootstrap intervals. This historical diagnostic is not part of the new gate cohort.}
\label{fig:supp_local_gate_controls}
\end{figure}

\subsection{Balanced extension of the matched credit-rule comparison}
\label{sec:credit_rule_extension}

After examining the released 1,024-update outcomes, we froze and committed a longer-budget protocol before running any extended trajectories. This is a protocol-led extension of previously observed seed blocks, not a fresh confirmatory cohort. All 720 algebraic trajectories were retained: twenty seed blocks, three task families, four delivery rules, Adam and SGD, and the union of the original development-selected and common rates. Matching and quartic targets retain the same balanced tree, initialization and example streams; nested targets retain their separately compatible trees. No rates were retuned.

Every trajectory was replayed from initialization and run for exactly 16,384 updates, preserving examples, minibatches, local rules, gradient clipping and coefficient bounds $[-2,2]$. All original checkpoint arrays, numerical curves and diagnostics through 1,024 updates reproduced. Additional observations occurred every 256 updates after 1,024. Terminal outcomes at 1,024, 4,096, 8,192 and 16,384 updates preserve the original endpoint definition. Secondary states minimized validation NMSE over the declared observations within each budget, with ties assigned to the earlier checkpoint. Test and full-domain population errors were used only for reporting. Independent recomputation of all 5,760 budget-outcome rows from saved weights reproduced test, validation and population NMSE exactly; all 2,880 validation-only selections were checked.

All twenty selected-rate Adam exact quartic fits reached test NMSE 0.0443--0.0484 by 4,096 updates, resolving the three slow outcomes at 1,024. At 16,384, the exact quartic mean was 0.0464 and its median 0.0461, compared with 0.9994 and 1.0769 for the initial profile (Supplementary Table~\ref{tab:credit_rule_extension_endpoints}). The original pairwise mean ordering was driven by slow exact seeds despite exact credit winning in fifteen of twenty selected-rate pairs. By 4,096, exact pairwise mean NMSE was 0.0232 versus 0.0240 for the initial profile, with eighteen of twenty exact wins. Common-rate pairwise differences remained small.

The calibrated-minus-exact quartet-minus-pairwise interaction at 16,384 was 0.9516 (95\% paired interval, 0.8829--1.0123) at terminal selected-rate states and 0.9028 (0.8371--0.9579) at terminal common-rate states. Validation-selected states gave smaller but still positive interactions: 0.7903 (0.6915--0.8797) and 0.7952 (0.7018--0.8775), respectively. All four contrasts were positive in every seed. The intervals use 10,000 paired whole-seed bootstrap draws and are descriptive; the seeds and their original outcomes had already been examined.

The comparisons remain constrained by time, parameter bounds and historical rates. Every selected-rate Adam initial-profile trajectory had contacted a coefficient bound by 16,384, whereas none of the corresponding exact trajectories had. Validation-selected initial-profile states still had substantial quartic and nested error, and the full record retains clipping and bound events. Extended exact-trained path fields remained almost one-dimensional for pairwise targets but multidirectional for quartic and nested targets; the original main-figure capture analysis is retained unchanged. These results strengthen the fixed-budget contrast and explain the slow exact runs, without proving that a bounded fixed profile fails under every optimizer or that full compartment resolution is uniquely necessary.

\begin{table}[tbp]
\centering
\small
\caption{Terminal outcomes in the balanced longer-budget follow-up. Rates are the original development-selected Adam rates. NMSE values are mean (median) across the same twenty previously examined seed blocks. The final column counts exact/profile fits with clean full-domain population NMSE below 0.02; this descriptive threshold was not used for stopping, selection or inference. Complete seedwise outcomes, validation-selected states, common-rate and SGD controls remain in Source Data.}
\label{tab:credit_rule_extension_endpoints}
\begin{tabular}{rlccc}
\toprule
Updates & Task & Exact NMSE & Initial-profile NMSE & Clean error $<0.02$ \\
 & & mean (median) & mean (median) & exact/profile counts \\
\midrule
1,024 & Pairwise & 0.0317 (0.0231) & 0.0236 (0.0236) & 19/20; 20/20 \\
1,024 & Quartic & 0.1516 (0.0465) & 0.9813 (1.0277) & 17/20; 0/20 \\
1,024 & Nested & 0.0230 (0.0229) & 0.7690 (0.7829) & 20/20; 0/20 \\
4,096 & Pairwise & 0.0232 (0.0231) & 0.0240 (0.0241) & 20/20; 20/20 \\
4,096 & Quartic & 0.0461 (0.0458) & 0.9624 (1.0205) & 20/20; 0/20 \\
4,096 & Nested & 0.0233 (0.0232) & 0.7913 (0.7853) & 20/20; 0/20 \\
8,192 & Pairwise & 0.0231 (0.0231) & 0.0243 (0.0243) & 20/20; 20/20 \\
8,192 & Quartic & 0.0463 (0.0463) & 0.9811 (1.0353) & 20/20; 0/20 \\
8,192 & Nested & 0.0233 (0.0233) & 0.8039 (0.7918) & 20/20; 0/20 \\
16,384 & Pairwise & 0.0232 (0.0232) & 0.0246 (0.0248) & 20/20; 20/20 \\
16,384 & Quartic & 0.0464 (0.0461) & 0.9994 (1.0769) & 20/20; 0/20 \\
16,384 & Nested & 0.0232 (0.0232) & 0.7944 (0.7826) & 20/20; 0/20 \\
\bottomrule
\end{tabular}
\end{table}

\begin{figure}[tbp]
\centering
\includegraphics[width=\textwidth]{supplementary/figure_S54_credit_budget_extension.pdf}
\caption{\textbf{Fixed profiles and exact credit under longer matched training budgets.} All 720 algebraic trajectories were replayed and extended to 16,384 updates, retaining twenty seed blocks, three tasks, four rules, two optimizers and the union of the original selected/common rates. \textbf{A--D}, Mean test NMSE with descriptive 95\% whole-seed bootstrap intervals for pairwise/quartic tasks under Adam at selected or common rates. Initial-sign controls use dotted curves. Vertical lines mark 1,024 updates; horizontal lines mark noise-only NMSE. \textbf{E}, Quartic-minus-pairwise difference in calibrated-minus-exact NMSE at four budgets. Circles use terminal states; squares use validation-selected states. \textbf{F}, Every paired quartic endpoint under exact or calibrated-profile credit at 1,024 and 16,384 updates; the diagonal marks unchanged error. Seeds were previously observed, rates remain fixed and coefficients remain bounded.}
\label{fig:supp_credit_rule_extension}
\end{figure}

\begin{figure}[tbp]
\centering
\includegraphics[width=\textwidth]{supplementary/figure_S55_credit_budget_controls.pdf}
\caption{\textbf{Nested targets and SGD controls in the balanced budget extension.} \textbf{A}, Nested targets at selected Adam rates. \textbf{B}, Pairwise targets at the common SGD rate. \textbf{C,D}, Quartic targets at selected and common SGD rates. The same previously observed twenty seeds, fixed rules and rate choices are retained. Curves show mean test NMSE from update 64 onward; bands give descriptive 95\% whole-seed bootstrap intervals. Vertical lines mark the original 1,024-update cap and horizontal lines the relevant noise floor. Panels state the task, optimizer and rate view; no outcome is excluded or used to retune rates. These controls retain optimizer dependence and bounded-parameter scope.}
\label{fig:supp_credit_rule_extension_controls}
\end{figure}

\subsection{Sensitivity of the measured ancestry-response test}
\label{note:measured_alignment_power}
We estimated detection sensitivity conditional on the observed sampling, using a Gaussian response model frozen before simulation. The analysis retained 125 partner observations of 102 presynaptic roots across thirteen scans and seven targets, including twenty target--partner combinations repeated across scans and two roots shared between targets. Each scan contained 136 repeated stimuli: 130 presented twice and six ten times. Original condition hashes identified 316 distinct repeated stimuli across scans. Pair tables, contact geometry, recording identities and the required observed responses are included in Source Data.

For each target arbor, an ancestry covariance was defined by the shared soma-to-common-ancestor path length divided by the geometric mean of the two soma-to-contact lengths. This is a positive-semidefinite Brownian path kernel; zero-depth contacts contributed independent diagonal components. The seven kernels were embedded in common presynaptic-root coordinates, summed and diagonal-normalized, so each shared root had one consistent latent response. Reliable tuning mixed independent Gaussian unit-specific and ancestry-correlated components with variance fractions $1-\lambda$ and $\lambda$, respectively. We tested $\lambda=0,0.05,\ldots,1$, retaining the actual partner sets and stimulus identities.

Measurement noise was calibrated to each recording's odd/even split-half Spearman reliability. Values were clipped to $[0,1]$ and converted to the Gaussian Pearson parameter $r_p=2\sin(\pi r_s/6)$; the two negative estimates remained archived and contributed zero reliable variance. With $h$ the mean inverse half-repeat count, trial responses had reliable variance $r_p$ and noise variance $(1-r_p)/h$. Condition-mean noise was divided by the actual repeat count. This matches the pooled Pearson second moment; its Spearman calibration is approximate because repeat counts vary. An independent 1,000-dataset audit found mean and maximum absolute Spearman discrepancies of 0.0020 and 0.0069, without retuning. A prespecified perfect-reliability comparison removed measurement noise.

In 4,000 simulated complete datasets, all pairwise similarities were computed from joint response matrices. The original partial-rank statistic controlled ranked Euclidean separation and contact-depth difference; scan effects were averaged within target. Positive detection required a positive mean effect and an exact two-sided signed-rank $P\leq0.05$ over the seven targets. All 128 sign assignments were enumerated. For seven nonzero untied effects, the minimum two-sided $P$ is 0.015625, and six of 128 null sign assignments reject at 0.05. This discreteness does not imply a universal power ceiling below one. Gaussian draws were paired across signal fractions and reliability scenarios; Wilson intervals describe Monte Carlo uncertainty.

With measured reliability, the first sustained tested 80\% crossing was $\lambda=0.60$, bracketed by 0.55--0.60. Descriptive interpolation gave $\lambda=0.589$ and a mean simulated target-level partial-rank effect of 0.249. Detection at $\lambda=0.60$ was 80.85\% (95\% Monte Carlo interval, 79.60--82.04\%); the lower interval first exceeded 80\% at $\lambda=0.65$. At maximal specified signal, detection was 96.23\% (95.59--96.77\%). The zero-signal two-sided rejection rate was 4.95\% (4.32--5.67\%). Perfect reliability reduced the interpolated variance fraction to 0.206, with a similar mean partial-rank effect of 0.248 at detection. These sensitivity estimates assume the specified Gaussian ancestry covariance, reliability calibration and known shared-input dependence; they do not exclude weaker effects, other covariance patterns, unknown biological dependence or alignment involving unobserved partners.

\begin{figure}[p]
\centering
\includegraphics[width=\textwidth]{supplementary/figure_S56_measured_alignment_power.pdf}
\caption{\textbf{Conditional sensitivity of the measured ancestry-response test.}
\textbf{A}, Detection of positive alignment as the fraction of reliable tuning variance carried by a shared-path covariance increases. Every simulated dataset preserves the thirteen scans, seven targets, observed partners, stimulus identities and repeats. Green uses measured repeat reliabilities; gray removes measurement noise.
\textbf{B}, The same detection probabilities versus the mean simulated seven-target partial-rank effect. This is a model-conditional effect axis, not a universal detectable correlation. The dashed line in \textbf{A,B} marks 80\%; bands are Wilson 95\% Monte Carlo intervals from 4,000 paired complete datasets. Detection requires an exact two-sided seven-target signed-rank $P\leq0.05$ and a positive mean effect.
\textbf{C}, Mean split-half Spearman reliability in 1,000 independent calibration datasets versus the observed value, for 125 partner records. The two negative estimates are assigned zero reliable variance (orange squares); the dotted segment shows that clipping. The identity line is not fitted. All partner-pair correlations come from joint response matrices; pairs are not inferential replicates. See Supplementary Note~\ref{note:measured_alignment_power} for the injected covariance and calibration assumptions.}
\label{fig:supp_measured_alignment_power}
\end{figure}

\subsection{Driving-force cancellation and the ancestry dependence of passive fields}
\label{sec:passive_field_diagnostics}
We reanalyzed the eight original cable operators used for the common-broadcast comparison. For each focal site, the perturbation was the original dose 0.001 with somatic voltage restored. The finite log-gradient difference divided by dose separates into an adjoint term and a driving-force term. Energy uses the original excitatory-area site weights and sums over focal perturbations. The full focal ancestry partition is defined separately for each shunt and is distinct from the seven selected routes in the budgeted anatomical dictionary.

Driving-force energy was 0.990--3.496 times full-field energy across cells; these are ratios of squared magnitudes, not amplitudes. Weighted centered correlations with the adjoint term were $-0.979$ to $-0.913$. The negative cross term explains why large components yield a smaller full response. The full response retained 0.9711--0.9993 of its energy in the focal partition. Reconstruction agrees with the original operators to $10^{-8}$; numerical decomposition closes to rounding error. Cellwise energies, correlations and partition captures are supplied in Source Data.

The passive corollary depends on a tree Green's function, a somatic adjoint source and a diagonal conductance change. Choosing other passive perturbations under the same assumptions does not make the anatomical field ancestry-neutral. The budgeted dictionary comparison measures compression within this response family. In the disjoint 47-cell cohort, actual ancestry exceeds the mean degree--depth surrogate in 38 cells, but in eleven cells more than half of 200 surrogates match or exceed the actual tree. This descriptive count uses tolerance $10^{-12}$ and retains the full within-cell distribution. A different error source, active dynamics or measured fields would be required to test beyond the present structural dependence.
